\documentclass[10pt,a4paper]{amsart}
\usepackage[margin=19mm]{geometry}
\usepackage{amsmath,amssymb,mathtools}
\usepackage[scr=rsfs,cal=euler]{mathalfa}
\usepackage{xfrac}
\usepackage[unicode,hidelinks,bookmarks=false]{hyperref}
\newcommand{\ie}{i.e.\ }
\newcommand{\cf}{cf.\ }
\newcommand{\resp}{respectively\ }
\newcommand{\Subsection}{\subsection}
\providecommand{\nobreakdash}{\nobreak\hbox{-}}
\setlength{\emergencystretch}{2em}
\multlinegap0pt
\usepackage{amssymb}
\usepackage{bm}
\usepackage{mathtools}
\usepackage{algorithm}\usepackage{algpseudocode}
\usepackage{xcolor}
\usepackage[shortlabels]{enumitem}
\usepackage{float}
\newtheorem{theorem}{Theorem}
\title{Iterative Methods for Computing the Moore--Penrose Inverse of Split-Quaternion Matrices with Applications}
\author{Salman Ahmadi-Asl, Valentin Leplat, Mohammad S. Alkousa}
\date{Source: arXiv:2607.29270v1 (CC BY 4.0)}
\begin{document}
\maketitle

\noindent\emph{Source and licence.} Iterative Methods for Computing the Moore--Penrose Inverse of Split-Quaternion Matrices with Applications. Version arXiv:2607.29270v1 (CC BY 4.0), distributed by arXiv under the Creative Commons Attribution 4.0 International licence, http://creativecommons.org/licenses/by/4.0/, as recorded in the arXiv licence metadata for this exact version and shown on the abstract page https://arxiv.org/abs/2607.29270v1. Full licence notice: \texttt{licenses/arxiv-2607.29270-CC-BY-4.0.txt}.\par\smallskip

\noindent\textit{Editorial scope.} Counted unit: the article's theorem theorem1 (source0.tex lines 640--675). The article's complete original proof of it is delivered with it (source0.tex lines 677--739, one \texttt{proof} environment of 63 lines). No mathematics is written, repaired or re-derived: the statement and the proof are the article's own lines, in the article's own notation, and no retained line is substituted or re-flowed.  Retained uncounted as the source-local context the proof needs: lines 626--628.  Nothing was omitted from any retained span, so the package carries no omission record.  No retained line contains a citation, so the wrapper carries no bibliography and no bibliography entry.  No retained line contains a figure environment, an image-inclusion command or drawing code, so no image is delivered and the package declares no asset dependency.  Editorial formatting only, in addition: an AMS \texttt{amsart} preamble that restates the article's own declarations and macros used by the retained lines, namely the environments \texttt{theorem} on separate counters, together with the standard packages those lines need.  The retained environments print in this document as Theorem 1 in order of appearance, whereas the article prints the same items with the numbering of its own full document.  The complete native source of the article as distributed in the arXiv e-print for this exact version is bundled unchanged as \texttt{sources/206468/source0.tex}.
\begin{equation} \label{eq:NS_real}
    X_{k+1} = X_k(2I_p-MX_k) = (2I_q-X_kM)X_k = 2X_k-X_kMX_k, \qquad X_0=\alpha M^T,
\end{equation}

\begin{theorem}[Exact Newton--Schulz convergence] \label{theorem1}
Let \(M\in\mathbb{R}^{p\times q}\) be nonzero and have rank \(r\), with positive singular values
\[
    \sigma_1\geq\sigma_2\geq\cdots\geq\sigma_r>0.
\]
Let \(\{X_k\}\) be generated by \eqref{eq:NS_real}. If
\begin{equation} \label{eq:alpha_interval}
    0<\alpha<\frac{2}{\sigma_1^2},
\end{equation}
then \(X_k\to M^\dagger\). More precisely, if
\[
    \theta_i:=1-\alpha\sigma_i^2, \qquad i=1,2,\ldots,r,
\]
then
\begin{equation}\label{eq:Xk_exact_formula}
    X_k =V_r\operatorname{diag}\!\left( \frac{1-\theta_1^{2^k}}{\sigma_1},\frac{1-\theta_2^{2^k}}{\sigma_2},\ldots, \frac{1-\theta_r^{2^k}}{\sigma_r} \right)U_r^T,
\end{equation}
where \(M=U_r\Sigma_rV_r^T\) is a thin SVD. In particular, with
\[
    P_L:=MM^\dagger=U_rU_r^T, \qquad P_R:=M^\dagger M=V_rV_r^T,
\]
one has
\begin{align}
    P_L-MX_k &=U_r\operatorname{diag}(\theta_1^{2^k},\theta_2^{2^k},\ldots,\theta_r^{2^k})U_r^T, \label{eq:left_projector_residual}
    \\ P_R-X_kM &=V_r\operatorname{diag}(\theta_1^{2^k},\theta_2^{2^k},\ldots,\theta_r^{2^k})V_r^T. \label{eq:right_projector_residual}
\end{align}
Hence, if
\[
    \rho_0:=\max_{1\leq i\leq r}|1-\alpha\sigma_i^2|<1,
\]
then
\begin{equation} \label{eq:projector_residual_rate}
    \|P_L-MX_k\|_2 =\|P_R-X_kM\|_2 =\rho_0^{2^k}.
\end{equation}
For the zero matrix, \(M^\dagger=0\) and the initialization gives \(X_k=0\) for all \(k\).
\end{theorem}

\begin{proof}
Let
\[
    M=U_r\Sigma_rV_r^T, \qquad \Sigma_r=\operatorname{diag}(\sigma_1,\sigma_2,\ldots,\sigma_r),
\]
be a thin SVD. Since
\[
    X_0=\alpha M^T =V_r(\alpha\Sigma_r)U_r^T,
\]
we first show by induction that every iterate has the form
\begin{equation} \label{eq:Xk_invariant_form}
    X_k=V_rD_kU_r^T,
\end{equation}
where \(D_k\) is diagonal. The statement is true at \(k=0\). If it holds at iteration \(k\), then
\begin{align*}
    X_{k+1}
    & =2V_rD_kU_r^T -V_rD_kU_r^T U_r\Sigma_rV_r^T V_rD_kU_r^T
    \\& =V_r\bigl(2D_k-D_k\Sigma_rD_k\bigr)U_r^T.
\end{align*}
Thus \eqref{eq:Xk_invariant_form} is preserved, with
\[
    D_{k+1}=2D_k-D_k\Sigma_rD_k.
\]
Writing \(D_k=\operatorname{diag}(d_{1,k},\ldots,d_{r,k})\), each singular direction evolves independently according to
\begin{equation} \label{eq:scalar_NS}
    d_{i,k+1}=d_{i,k}(2-\sigma_i d_{i,k}), \qquad d_{i,0}=\alpha\sigma_i.
\end{equation}
Define
\[
    e_{i,k}:=1-\sigma_i d_{i,k}.
\]
Then \eqref{eq:scalar_NS} gives
\begin{align*}
    e_{i,k+1}
    &=1-\sigma_i d_{i,k}(2-\sigma_i d_{i,k})\\
    &=1-2\sigma_i d_{i,k}+\sigma_i^2d_{i,k}^2\\
    &=(1-\sigma_i d_{i,k})^2 =e_{i,k}^2.
\end{align*}
Since \(e_{i,0}=1-\alpha\sigma_i^2=\theta_i\), repeated squaring yields
\[
    e_{i,k}=\theta_i^{2^k}.
\]
Solving for \(d_{i,k}\) gives
\[
    d_{i,k}=\frac{1-\theta_i^{2^k}}{\sigma_i},
\]
which proves \eqref{eq:Xk_exact_formula}.

Condition \eqref{eq:alpha_interval} implies
\[
    -1<1-\alpha\sigma_1^2 \leq 1-\alpha\sigma_i^2<1
\]
for each \(i\), so \(|\theta_i|<1\). Therefore \(\theta_i^{2^k}\to0\), and \eqref{eq:Xk_exact_formula} gives
\[
    X_k\longrightarrow V_r\Sigma_r^{-1}U_r^T =M^\dagger.
\]
Finally,
\begin{align*}
    MX_k &=U_r\Sigma_rD_kU_r^T =U_r\operatorname{diag}(1-\theta_i^{2^k})U_r^T,
    \\ X_kM &=V_rD_k\Sigma_rV_r^T =V_r\operatorname{diag}(1-\theta_i^{2^k})V_r^T,
\end{align*}
which proves \eqref{eq:left_projector_residual}--\eqref{eq:projector_residual_rate}.
\end{proof}

\end{document}
